\chapter{RESULTS AND DISCUSSION}

\section{System Interface Showcase}
The system was successfully deployed and verified. The primary user interfaces are detailed below:
\begin{itemize}
    \item \textbf{Home Page Landing:} Renders a premium dark-mode dashboard showcasing core performance KPIs: 98.4\% AI Accuracy, $<$16ms latency, 13 tracked exercises, and 60+ FPS rendering speeds. Interactive animations introduce the system features.
    \item \textbf{Live Workout Playground:} Renders a zero-scroll dashboard splitting the screen between the live webcam stream (with real-time skeleton joints and form alert text) and the workout HUD control panel. The control panel allows adjusting target reps, displays active exercise details, and runs a stopwatch session timer.
    \item \textbf{Video Analysis Studio:} Features an intuitive drag-and-drop file dashboard that processes pre-recorded MP4 clips, providing skeleton overlays and rep counting for video file reviews.
    \item \textbf{Deep Learning Models Page:} Visualizes the custom sequence classification pipeline. It details how joint keypoint coordinate arrays are compiled over 30-frame sequence lengths and processed by the 1D-CNN + LSTM architecture.
    \item \textbf{Architecture Page:} Explains the technical data flow path, listing the components of the full-stack system and tracing the path from raw pixel camera streams to metric evaluations.
\end{itemize}

\section{Sequence Classification Model Performance}
The sequence classification model was trained on a dataset of over 10,300 sequence samples spanning 13 different exercise movements. Each sample is a sequence of 30 frames containing 132 keypoint coordinates.
The training hyperparameters were:
\begin{itemize}
    \item \textbf{Optimizer:} Adam ($\text{learning rate} = 0.001$)
    \item \textbf{Loss Function:} Categorical Cross-Entropy
    \item \textbf{Batch Size:} 64
    \item \textbf{Epochs:} 50 (with early stopping patience of 7)
\end{itemize}
The model achieved a training accuracy of \textbf{99.1\%} and a test validation accuracy of \textbf{98.4\%}. The confusion matrix indicated high precision across similar movement signatures (e.g. bicep curl vs shoulder press), thanks to the temporal features extracted by the LSTM layer.

\section{System Performance and Benchmarks}
Performance evaluations were conducted comparing CPU-only execution and GPU-accelerated execution. Table 5.1 details the metrics:

\begin{table}[h]
\centering
\begin{tabular}{|p{4.5cm}|p{5.5cm}|p{5.5cm}|}
\hline
\textbf{Metric} & \textbf{CPU (Intel Core i5)} & \textbf{GPU (NVIDIA CUDA)} \\ \hline
Pose Estimation Time & $\sim 12 - 18$ ms / frame & $\sim 3 - 5$ ms / frame \\ \hline
Model Inference Time & $\sim 14 - 22$ ms / frame & $\sim 2 - 4$ ms / frame \\ \hline
Total Latency per Frame & $\sim 26 - 40$ ms & $\sim 5 - 9$ ms \\ \hline
Stream Frame Rate (FPS) & $\sim 25 - 30$ FPS & $\sim 60+$ FPS \\ \hline
Memory Utilization & $\sim 350$ MB RAM & $\sim 120$ MB VRAM \\ \hline
\end{tabular}
\caption{System Inference Performance Benchmarks.}
\label{tab:performance_benchmarks}
\end{table}

The benchmarking indicates that although CPU execution is sufficient to maintain standard camera frame rates (30 FPS), GPU acceleration provides a fluid, high-performance experience ($>60$ FPS) with minimal latency.

\section{Discussion of Key Findings}
Several critical architectural insights were gained during system testing:
\begin{itemize}
    \item \textbf{Hybrid Edge-Cloud Scalability:} Standard computer vision applications struggle with scale due to server-side GPU costs. By utilizing lightweight local client interfaces linked to an edge backend, the system eliminates network latency. Pushing drawn shapes directly on the frame in the backend thread guarantees that overlay coordinates match joint lines, preventing drift.
    \item \textbf{Sequence Length Optimization:} Setting the temporal window to 30 frames ($1.0$ second of video at 30 FPS) balances classification speed and accuracy. Shorter sequences (e.g., 15 frames) resulted in higher false triggers, whereas longer sequences (e.g., 60 frames) added classification delay, preventing real-time exercise adaptation.
    \item \textbf{Trigonometry vs Deep Learning:} Relying on hybrid combinations is highly effective. Biomechanical angles are computed using pure mathematical formulas, ensuring mathematical consistency. Deep learning classifies the general activity type, allowing the system to dynamically toggle the active state-machine mathematical thresholds.
\end{itemize}
\newpage
